\documentclass[11pt]{article}
\usepackage[margin=1in]{geometry}
\usepackage{amsmath,amssymb,bm}
\usepackage{booktabs}
\usepackage[colorlinks=true,linkcolor=blue,citecolor=blue]{hyperref}
\newcommand{\D}{\mathrm{D}}
\newcommand{\dd}{\partial}
\newcommand{\er}{\bm{e}_r}
\newcommand{\etheta}{\bm{e}_\theta}
\newcommand{\ez}{\bm{e}_z}
\newcommand{\om}{\bm{\omega}}
\newcommand{\ub}{\bm{u}}
\title{Two theoretical limits of the underexpanded-jet campaign:\\
(i) the NPR lower bound, and (ii) why the axisymmetric (RZ) model\\
is valid only in the near field --- a vorticity-equation argument}
\author{Cerisse U\_jet campaign notes}
\date{2 July 2026}
\begin{document}
\maketitle

%=====================================================================
\section{The lower bound of NPR: the choking threshold}

For a convergent (sonic) nozzle discharging into quiescent ambient at
$p_\infty$, the jet is \emph{underexpanded} only if the nozzle is
choked. Isentropic flow gives the exit static pressure of a choked
nozzle, $p_e/p_0 = \left(\tfrac{2}{\gamma+1}\right)^{\gamma/(\gamma-1)}$,
so the exit is exactly pressure-matched ($p_e = p_\infty$) when
\begin{equation}
  \mathrm{NPR} \equiv \frac{p_0}{p_\infty}
  = \left(\frac{\gamma+1}{2}\right)^{\frac{\gamma}{\gamma-1}}
  = 1.892\,93 \qquad (\gamma = 1.4).
  \label{eq:choke}
\end{equation}
This is the ``$1.8\mathrm{xx}$'' bound quoted in the literature
(e.g.\ Franquet \emph{et al.}, \emph{Prog.\ Aerosp.\ Sci.} 2015):
\begin{itemize}
\item $\mathrm{NPR} < 1.893$: nozzle unchoked, exit subsonic, ordinary
  subsonic jet --- no shock-cell train exists at all;
\item $\mathrm{NPR} = 1.893$: perfectly expanded sonic jet
  ($p_e/p_\infty = 1$, $M_j = 1$);
\item $\mathrm{NPR} > 1.893$: underexpanded sonic jet, with the
  fully-expanded Mach number
  $M_j = \sqrt{\tfrac{2}{\gamma-1}\!\left[\mathrm{NPR}^{(\gamma-1)/\gamma}-1\right]}$.
\end{itemize}

Our campaign uses $\mathrm{NPR}=p_0/p_\infty$ directly, so the current
case matrix sits as follows:
\begin{center}
\begin{tabular}{lccl}
\toprule
NPR & $p_e/p_\infty$ & $M_j$ & regime \\
\midrule
1.893 & 1.000 & 1.000 & threshold (perfectly expanded) \\
2     & 1.057 & 1.046 & \textbf{+5.7\% above the limit} --- weakest possible wave train \\
3     & 1.585 & 1.358 & moderately underexpanded, regular reflection \\
3.8   & 2.0   & 1.53  & empirical Mach-disk onset ($p_e/p_\infty \simeq 2$) \\
4--5  & 2.1--2.6 & 1.56--1.71 & Mach disk established \\
\bottomrule
\end{tabular}
\end{center}
Thus \textbf{NPR $=2$ is essentially at the lower boundary of the
underexpanded regime}: the exit over-pressure is only $5.7\%$, the
wave train consists of weak isentropic compression/expansion waves in
regular reflection (no embedded shocks strong enough for a Mach disk),
and all density gradients are smooth --- which is also why
gradient-based AMR tagging self-limits for this case.

%=====================================================================
\section{Why RZ-axisymmetry is a near-field model: the vorticity equation}

\subsection{Setup}
Compressible vorticity transport (curl of the momentum equation):
\begin{equation}
  \frac{\D\om}{\D t}
  \;=\;
  \underbrace{(\om\!\cdot\!\nabla)\ub}_{\text{stretching/tilting}}
  \;-\;
  \underbrace{\om(\nabla\!\cdot\!\ub)}_{\text{dilatation}}
  \;+\;
  \underbrace{\frac{\nabla\rho\times\nabla p}{\rho^{2}}}_{\text{baroclinic}}
  \;+\;
  \underbrace{\nabla\times\!\Big(\tfrac{1}{\rho}\nabla\!\cdot\!\bm{\tau}\Big)}_{\text{viscous}} .
  \label{eq:vort3d}
\end{equation}
In cylindrical coordinates with the axisymmetric, swirl-free ansatz of
the RZ solver ($\dd_\theta \equiv 0$, $u_\theta \equiv 0$):
\begin{equation}
  \ub = u_r\,\er + u_z\,\ez,
  \qquad
  \om = \nabla\times\ub = \omega_\theta\,\etheta,
  \qquad
  \omega_\theta = \frac{\dd u_r}{\dd z}-\frac{\dd u_z}{\dd r}.
\end{equation}
\emph{All vorticity is azimuthal}: vortex lines are perfect circles
around the axis (ring vortices), by construction and forever.

\subsection{The stretching term collapses to ring dilation}
Evaluate the 3-D production term under the ansatz. Since the scalar
components are $\theta$-independent but $\dd_\theta\er=\etheta$,
\begin{equation}
  (\om\!\cdot\!\nabla)\ub
  = \frac{\omega_\theta}{r}\,\dd_\theta\!\left(u_r\er+u_z\ez\right)
  = \frac{u_r\,\omega_\theta}{r}\,\etheta .
  \label{eq:stretchRZ}
\end{equation}
The tensorial stretching $(\om\!\cdot\!\nabla)\ub = \mathsf{S}\om + \tfrac12\om\times\om\dots$
degenerates into the single scalar $u_r/r$: the only way an
axisymmetric vortex ring can ``stretch'' is by increasing its own
radius. Substituting \eqref{eq:stretchRZ} into \eqref{eq:vort3d} and
using continuity $\D\rho/\D t=-\rho\nabla\!\cdot\!\ub$ together with
$\D r/\D t = u_r$ gives the scalar transport law
\begin{equation}
  \boxed{\;
  \frac{\D}{\D t}\!\left(\frac{\omega_\theta}{\rho\, r}\right)
  = \frac{1}{\rho^{3} r}\left(\nabla\rho\times\nabla p\right)_\theta
  + \frac{1}{\rho r}\Big[\nabla\times\big(\tfrac1\rho\nabla\!\cdot\!\bm\tau\big)\Big]_\theta
  \;}
  \label{eq:invariant}
\end{equation}
i.e.\ \textbf{for inviscid, barotropic flow $\omega_\theta/(\rho r)$ is a
material invariant} (the compressible generalisation of the classical
axisymmetric invariant $\omega_\theta/r$).

\subsection{Consequences: no amplification, no cascade}
\paragraph{(a) Algebraic bound on vorticity.}
Integrating \eqref{eq:invariant} along a particle path,
\begin{equation}
  \omega_\theta(t)
  = \omega_\theta(0)\,\frac{\rho(t)\,r(t)}{\rho(0)\,r(0)}
  \;+\; \text{(baroclinic integral)} .
\end{equation}
Both $\rho$ and $r$ are bounded, so vorticity can grow only
\emph{algebraically}, through ring dilation or compression. Contrast
3-D: writing $\om=\omega\,\bm e_\omega$, the stretching term gives
$\D\omega/\D t = (\bm e_\omega\!\cdot\!\mathsf{S}\,\bm e_\omega)\,\omega + \dots$,
and the well-known preferential alignment of $\bm e_\omega$ with the
extensional strain eigenvector makes the growth rate positive on
average --- \emph{exponential} amplification, terminated only at the
Kolmogorov scale. This is the engine of the direct energy cascade; RZ
removes it identically.

\paragraph{(b) Enstrophy budget.}
The inviscid enstrophy production in 3-D is
$\int \om\,\mathsf{S}\,\om\,\mathrm dV$, positive on average
($\propto -\langle S^3\rangle$, velocity-derivative skewness
$\simeq -0.4$ in real turbulence). Under the RZ ansatz
\begin{equation}
  \om\,\mathsf{S}\,\om = \omega_\theta^{2}\,S_{\theta\theta}
  = \omega_\theta^{2}\,\frac{u_r}{r},
\end{equation}
which is sign-indefinite and tied to the \emph{mean} radial motion ---
there is no self-amplification mechanism. Enstrophy is (up to
baroclinicity and viscosity) a Casimir-constrained quantity, exactly
as in 2-D turbulence.

\paragraph{(c) Wrong cascade direction.}
With stretching absent, the $(r,z)$-plane dynamics possess two
quadratic inviscid invariants (kinetic energy and the family
$\int \rho f(\omega_\theta/\rho r)\,\mathrm dV$). By the standard
Fj\o rtoft/Kraichnan argument this forces a \emph{dual cascade}:
energy flows to \emph{large} scales (vortex pairing and merging),
enstrophy to small scales. The phenomenology in our RZ solutions ---
shear-layer rings that pair, grow, and form large coherent ``puffing''
structures instead of breaking down --- is not a numerical artefact
but the \emph{correct} solution of the axisymmetric equations, which
are simply the wrong equations for a turbulent far field. A real jet
obeys the single-invariant 3-D Richardson--Kolmogorov direct cascade.

\paragraph{(d) Modal truncation.}
Decomposing any disturbance as $\sum_m \hat q(r,z)\,e^{im\theta}$, the
RZ solver retains only $m=0$. Linear stability theory of round jets
(Batchelor \& Gill 1962; Michalke 1971) shows: for the thin annular
shear layer near the lip ($\delta_\omega/R \ll 1$) the $m=0$ and
$|m|=1$ modes have comparable growth rates --- so the \emph{initial}
roll-up into rings is genuinely captured; but for the thick,
developed profile downstream the \emph{only} unstable mode is the
helical $|m|=1$, and the dominant secondary instabilities of the ring
street (Widnall-type azimuthal core waves, elliptic instability of
the cores, braid streamwise vortices) are all $m\neq0$. The RZ model
therefore \emph{cannot even represent linearly} the instabilities
that govern transition and far-field dynamics.

\subsection{Why the near field survives}
The same decomposition explains the region of validity:
\begin{enumerate}
\item \textbf{Wave-system gasdynamics is $m=0$ physics.} The
  Prandtl--Meyer fan, intercepting/barrel shock, Mach disk and
  reflected shock are steady solutions of the axisymmetric Euler
  equations (method of characteristics + Rankine--Hugoniot); the real
  jet's near-field wave train \emph{is} axisymmetric to leading order.
\item \textbf{Near-field vorticity is baroclinic, and its $m=0$ part
  is retained.} Equation \eqref{eq:invariant} keeps the full
  axisymmetric baroclinic source $(\nabla\rho\times\nabla p)_\theta$:
  vorticity generated at curved shocks and at the Mach-disk slip line
  is captured --- these set the shear layers that bound the first
  cells.
\item \textbf{Initial K--H roll-up is quasi-planar.} With
  $\delta_\omega \ll R$ the annular layer behaves locally like a plane
  mixing layer whose primary instability is two-dimensional
  (Michalke); rings are indeed observed in real jets for the first
  1--2 pairings.
\item \textbf{Failure begins at the first 3-D secondary instability.}
  After $\mathcal O(1)$ pairings the Widnall instability grows on the
  ring-turnover timescale; beyond that station stretching, azimuthal
  modes and the direct cascade --- all absent in RZ --- control the
  physics.
\end{enumerate}
Empirically, in this campaign (NPR $=3$): centerline and field data
agree between RZ and full 3-D to $<3\%$ up to $z\approx1.8\,D_e$
(through the first shock cell and Mach disk), while at $z=20\,D_e$
the RZ centerline velocity is $1.8\times$ the 3-D value and the
high-NPR RZ far field exhibits $25$--$40\%$ rms coherent puffing ---
precisely the inverse-cascade signature predicted above.

\paragraph{Summary.}
RZ axisymmetry deletes exactly the two ingredients that make jet
turbulence three-dimensional --- exponential vortex stretching
(replaced by the material invariant $\omega_\theta/\rho r$,
Eq.~\ref{eq:invariant}) and all $m\neq0$ azimuthal modes --- while
retaining everything that fixes the near-field wave structure
(axisymmetric gasdynamics and $m=0$ baroclinicity). It is therefore a
quantitatively reliable \emph{near-field} model (first few shock
cells) and a qualitatively wrong \emph{far-field} model, independent
of grid resolution.

\end{document}
